\summarychapter{Summary of certain properties of finite-dimensional Lie algebras over a field of characteristic 0}{Summary}{of certain properties of finite-dimensional Lie algebras over a field of characteristic 0.}


Let g be a Lie algebra, r its radical, n its largest nilpotent ideal, s its nilpotent radical and t the orthogonal of g relative to the Killing form. Then r, n, s, t are characteristic ideals and \(r \supset t \supset n \supset s\) .

\begin{enumerate}
\def\labelenumi{(\Roman{enumi})}
\tightlist
\item
  Any one of the following properties characterizes semi-simple Lie algebras:
\end{enumerate}

\begin{enumerate}
\def\labelenumi{(\arabic{enumi})}
\tightlist
\item
  \(r = \{0\}\) ; (2) \(n = \{0\}\) ; (3) \(t = \{0\}\) ; (4) every commutative ideal of g is zero; (5) the algebra g is isomorphic to a product of simple Lie algebras; (6) every finite-dimensional representation of g is semi-simple.
\end{enumerate}

\begin{enumerate}
\def\labelenumi{(\Roman{enumi})}
\setcounter{enumi}{1}
\tightlist
\item
  Any one of the following properties characterizes reductive Lie algebras:
\end{enumerate}

\begin{enumerate}
\def\labelenumi{(\arabic{enumi})}
\tightlist
\item
  \(\mathfrak{s} = \{0\}\) ; (2) \(\mathfrak{r}\) is the centre of \(\mathfrak{g}\) ; (3) \(\mathcal{D}\mathfrak{g}\) is semi-simple; (4) \(\mathfrak{g}\) is the product of a semi-simple algebra and a commutative algebra; (5) the adjoint representation of \(\mathfrak{g}\) is semi-simple; (6) \(\mathfrak{g}\) has a finite-dimensional representation such that the associated bilinear form is non-degenerate; (7) \(\mathfrak{g}\) has a finite-dimensional semi-simple faithful representation.
\end{enumerate}

\begin{enumerate}
\def\labelenumi{(\Roman{enumi})}
\setcounter{enumi}{2}
\tightlist
\item
  Any one of the following properties characterizes solvable Lie algebras:
\end{enumerate}

\begin{enumerate}
\def\labelenumi{(\arabic{enumi})}
\tightlist
\item
  \(\mathcal{D}^{p}\mathfrak{g} = \{0\}\) for sufficiently large p; (2) there exists a decreasing sequence \(g = g_{0} \supset g_{1} \supset \cdots \supset g_{n} = \{0\}\) of ideals of g such that the algebras \(g_{i}/g_{i+1}\) are commutative; (3) there exists a decreasing sequence
\end{enumerate}

\[
\mathfrak {g} = \mathfrak {g} _ {0} ^ {\prime} \supset \mathfrak {g} _ {1} ^ {\prime} \supset \dots \supset \mathfrak {g} _ {n ^ {\prime}} ^ {\prime} = \{0 \}
\]

of subalgebras of \(\mathfrak{g}\) such that \(\mathfrak{g}_{i+1}^{\prime}\) is an ideal of \(\mathfrak{g}_i^{\prime}\) and \(\mathfrak{g}_i^{\prime}/\mathfrak{g}_{i+1}^{\prime}\) is commutative; (4) there exists a decreasing sequence \(\mathfrak{g} = \mathfrak{g}_0'' \supset \mathfrak{g}_1'' \supset \cdots \supset \mathfrak{g}_{n''}'' = \{0\}\) of subalgebras of \(\mathfrak{g}\) such that \(\mathfrak{g}_{i+1}''\) is an ideal of codimension 1 in \(\mathfrak{g}_i''\) ; (5) \(\mathfrak{g} \supset \mathcal{D}\mathfrak{g}\) ; (6) \(\mathcal{D}\mathfrak{g}\) is nilpotent.

\begin{enumerate}
\def\labelenumi{(\Roman{enumi})}
\setcounter{enumi}{3}
\tightlist
\item
  Any one of the following properties characterizes nilpotent Lie algebras:
\end{enumerate}

\begin{enumerate}
\def\labelenumi{(\arabic{enumi})}
\tightlist
\item
  \(\mathcal{C}_{g}^{p} = \{0\}\) for sufficiently large p; (2) \(\mathcal{C}_{p}\mathfrak{g} = \mathfrak{g}\) for sufficiently large p; (3) there exists a decreasing sequence \(g = g_{0} \supset g_{1} \supset \cdots \supset g_{p} = \{0\}\) of ideals of g such that \([g, g_{i}] \subset g_{i+1}\) ; (4) there exists a decreasing sequence \(g = g_{0}' \supset g_{1}' \supset \cdots \supset g_{p'}' = \{0\}\) of ideals of g such that \([g, g_{i}'] \subset g_{i+1}'\) and the \(g_{i}'/g_{i+1}'\) are 1-dimensional; (5) there exists an integer i such that
\end{enumerate}

\[
\left(\operatorname{ad} x _ {1}\right) \circ \left(\operatorname{ad} x _ {2}\right) \circ \dots \circ \left(\operatorname{ad} x _ {i}\right) = 0
\]

for all \(x_{1},\ldots ,x_{i}\) in \(\mathfrak{g}\) ; (6) for all \(x\in \mathfrak{g}\) , ad \(x\) is nilpotent.

\begin{enumerate}
\def\labelenumi{(\Alph{enumi})}
\setcounter{enumi}{21}
\tightlist
\item
  g commutative \(\Rightarrow\) g nilpotent \(\Rightarrow\) the Killing form of g is zero \(\Rightarrow\) g solvable.
\end{enumerate}

g commutative \(\Rightarrow\) g reductive.

g semi-simple \(\Rightarrow\) g reductive.

\textbf{(VI) Characterizations of r:}

\begin{enumerate}
\def\labelenumi{(\arabic{enumi})}
\tightlist
\item
  r is the largest solvable ideal of g; (2) r is the smallest ideal such that g/r is semi-simple; (3) r is the only solvable ideal such that g/r is semi-simple; (4) r is the orthogonal of \(\mathcal{D}\mathfrak{g}\) relative to the Killing form.
\end{enumerate}

\textbf{(VII) Characterizations of n:}

\begin{enumerate}
\def\labelenumi{(\arabic{enumi})}
\tightlist
\item
  n is the largest nilpotent ideal of g; (2) n is the largest nilpotent ideal of r; (3) n is the set of \(x \in r\) such that \(\operatorname{ad}_{\mathfrak{g}} x\) is nilpotent; (4) n is the set of \(x \in r\) such that \(\operatorname{ad}_{\mathfrak{r}} x\) is nilpotent; (5) n is the largest ideal of g such that, for all \(x \in n\) , \(\operatorname{ad}_{\mathfrak{g}} x\) is nilpotent; (6) n is the set of \(x \in g\) such that \(\operatorname{ad}_{\mathfrak{g}} x\) belongs to the radical of the associative algebra generated by 1 and the \(\operatorname{ad}_{\mathfrak{g}} y\,(y \in g)\) .
\end{enumerate}

\textbf{(VIII) Characterizations of s:}

\begin{enumerate}
\def\labelenumi{(\arabic{enumi})}
\tightlist
\item
  \(\mathfrak{s}\) is the intersection of the kernels of the finite-dimensional simple representations of \(\mathfrak{g}\) ; (2) \(\mathfrak{s}\) is the smallest of the kernels of the finite-dimensional semi-simple representations of \(\mathfrak{g}\) ; (3) \(\mathfrak{s}\) is the intersection of the largest nilpotency ideals of the finite-dimensional representations of \(\mathfrak{g}\) ; (4) \(\mathfrak{s}\) is the smallest ideal of \(\mathfrak{g}\) such that \(\mathfrak{g} / \mathfrak{s}\) is reductive; (5) \(\mathfrak{s} = \mathfrak{r} \cap \mathcal{D}\mathfrak{g}\) ; (6) \(\mathfrak{s} = [\mathfrak{r}, \mathfrak{g}]\) ; (7) \(\mathfrak{s}\) is the intersection of the orthogonals of \(\mathfrak{g}\) relative to the bilinear forms associated with the finite-dimensional representations of \(\mathfrak{g}\) .
\end{enumerate}


